  let tex = fs.readFileSync('head.tex','utf8');

  tex += `



\documentclass[12pt,a4paper]{article}

\usepackage[a4paper,margin=1in]{geometry}
\usepackage[dvipsnames]{xcolor}
\usepackage{tcolorbox}
\usepackage{longtable}
\usepackage{booktabs}
\usepackage{array}
\usepackage{hyperref}
\usepackage{titlesec}
\usepackage{fancyhdr}

% ---------------- COLORS ----------------
\definecolor{mainblue}{RGB}{0,70,140}
\definecolor{lightblue}{RGB}{235,245,255}
\definecolor{darkgray}{RGB}{60,60,60}

% ---------------- STYLE ----------------
\titleformat{\section}
{\Large\bfseries\color{mainblue}}
{\thesection}{1em}{}

\titleformat{\subsection}
{\large\bfseries\color{darkgray}}
{\thesubsection}{1em}{}






\pagestyle{fancy}
\fancyhf{}
\rhead{Pharmacy V12}
\lhead{Termux System}
\cfoot{\thepage}

% ---------------- TITLE ----------------
\begin{document}

\begin{titlepage}
\centering

\vspace*{2cm}

{\Huge\bfseries\color{mainblue}
Pharmacy Management System V12\par}

\vspace{0.5cm}

{\Large Termux • Node.js • SAS • LaTeX Integration\par}

\vspace{1.5cm}

\begin{tcolorbox}[
colback=lightblue,
colframe=mainblue,
width=0.85\textwidth,
arc=4mm,
boxrule=0.8pt]

\centering
\textbf{System Design Document}\\
Inventory • Sales • Analytics • Reporting • Backup • Health Monitoring

\end{tcolorbox}

\vfill

{\large Prepared by: Babiker Osman\par}
{\large \today\par}

\end{titlepage}

% ---------------- EXECUTIVE SUMMARY ----------------
\section{Executive Summary}

\begin{tcolorbox}[
colback=lightblue,
colframe=mainblue,
title=Overview]

The Pharmacy Management System V12 is a lightweight, offline-ready pharmacy platform designed for Termux environments. It integrates inventory management, sales processing, analytics, SAS export, backups, and LaTeX reporting into a single unified system.

\end{tcolorbox}

% ---------------- SYSTEM ARCHITECTURE ----------------
\section{System Architecture}

The system consists of modular components:

\begin{itemize}
\item Inventory Engine
\item Sales Engine
\item Analytics Engine
\item Reporting Engine
\item Web Dashboard
\item SAS Export Engine
\item Backup System
\item Health Monitoring
\end{itemize}

% ---------------- MENU ----------------
\section{Operational Menu Structure}

\begin{longtable}{cl}
\toprule
\textbf{Option} & \textbf{Function} \\
\midrule

1  & Add Medicine \\
2  & List Medicines \\
3  & Sell Medicine \\
4  & Search Medicine \\
5  & Alerts \\
6  & Analytics \\
7  & Web Dashboard \\
8  & Reports (TXT) \\
9  & Stop Server \\
10 & Exit \\
11 & SAS Export \\
12 & Backup System \\
13 & System Health Check \\

\bottomrule
\end{longtable}

% ---------------- FEATURES ----------------
\section{System Highlights}

\begin{tcolorbox}[
colback=lightblue,
colframe=mainblue,
title=Key Features]

\begin{itemize}
\item Fully offline Termux-compatible system
\item Lightweight JSON-based storage
\item Web dashboard for real-time monitoring
\item SAS-compatible data export (CSV)
\item Automated backup system
\item Audit logging for all actions
\item Health and inventory monitoring
\item TXT + LaTeX reporting support
\item Modular and extendable architecture
\end{itemize}

\end{tcolorbox}

% ---------------- REPORTING ----------------
\section{Reporting System}

\subsection{TXT Reports}
Generated for quick operational summaries:
\begin{verbatim}
reports/pharmacy_report.txt
\end{verbatim}

\subsection{LaTeX Reports}
Professional PDF output using:
\begin{verbatim}
~/compiletex1.sh
\end{verbatim}

% ---------------- SAS ----------------
\section{SAS Integration}

Exports structured datasets for analytics:

\begin{verbatim}
exports/medicines_sas.csv
exports/sales_sas.csv
\end{verbatim}

Compatible with SAS, Python, R, and Excel.

% ---------------- BACKUP ----------------
\section{Backup System}

Creates timestamped backups of:

\begin{itemize}
\item Medicines database
\item Sales records
\item Audit logs
\end{itemize}

Ensures full data recovery capability.

% ---------------- HEALTH ----------------
\section{System Health Check}

Provides system diagnostics:

\begin{itemize}
\item Total medicines
\item Total sales
\item Revenue summary
\item Low stock alerts
\end{itemize}




\section*{Dashboard Design Philosophy}

The Pharmacy Management System V12 dashboard combines a healthcare-oriented visual identity with data analytics principles inspired by SAS reporting environments.

\subsection*{Healthcare Theme}

The primary visual theme is based on medical green tones, representing:

\begin{itemize}
\item Patient care
\item Healthcare services
\item Medication management
\item Safety and reliability
\item Pharmacy operations
\end{itemize}

The dashboard background employs a light green palette designed to reduce visual fatigue while maintaining a professional clinical appearance.

\subsection*{SAS Analytics Theme}

The secondary visual theme incorporates SAS-inspired blue accents to emphasize:

\begin{itemize}
\item Business intelligence
\item Statistical reporting
\item Data analytics
\item Revenue monitoring
\item Decision support systems
\end{itemize}

These elements reinforce the analytical capabilities of the platform and support management-oriented decision making.

\subsection*{Dashboard Components}

\begin{longtable}{ll}
\toprule
Component & Color Theme \\
\midrule
Medicine Inventory & Medical Green \\
Sales Activity & Teal \\
Revenue Analytics & SAS Blue \\
System Alerts & Red \\
\bottomrule
\end{longtable}

\subsection*{Design Objectives}

The combined Healthcare--SAS design philosophy aims to achieve:

\begin{itemize}
\item Rapid interpretation of key performance indicators
\item Improved user experience on mobile devices
\item Consistent pharmacy branding
\item Enhanced analytical visibility
\item Professional management reporting
\end{itemize}

This dual-theme approach creates a balanced environment that reflects both the operational requirements of a pharmacy and the analytical strengths of modern data-driven management systems.



\section{Python Simulation Models (Pharmacy V12) }

The Pharmacy Management System integrates advanced Python-based simulation models to support decision-making, forecasting, and operational optimization.

These models are executed from the Termux-based Node.js application and operate on real pharmacy data stored in JSON format.

\subsection{ Inventory Risk Simulation}

This model estimates the probability of stock-out events using Monte Carlo simulation.

\textbf{Purpose:}
\begin{itemize}
\item Measure risk of insufficient stock
\item Simulate demand variability
\item Identify vulnerable medicines
\end{itemize}

\textbf{Method:}
Random demand is generated over 1000 iterations and compared with current stock levels.

\subsection{ Demand Forecasting Model}

Uses statistical averaging and stochastic simulation to estimate future demand.

\textbf{Purpose:}
\begin{itemize}
\item Predict future pharmacy demand
\item Support procurement planning
\end{itemize}

\textbf{Method:}
Monte Carlo simulation around historical mean sales.

\subsection{ Reorder Point Model}

Calculates when a medicine should be reordered.

\textbf{Formula:}
\[
ROP = (Daily\ Demand \times Lead\ Time)
\]

\textbf{Purpose:}
\begin{itemize}
\item Prevent stock depletion
\item Optimize inventory replenishment
\end{itemize}

\subsection{ Safety Stock Model}

Determines buffer stock required to absorb demand fluctuations.

\textbf{Formula:}
\[
Safety\ Stock = Z \times \sqrt{Lead\ Time} \times Demand
\]

\textbf{Purpose:}
\begin{itemize}
\item Reduce stock-out risk
\item Maintain service continuity
\end{itemize}

\subsection{ Expiry Risk Simulation}

Analyzes medicine expiry timelines and flags high-risk items.

\textbf{Purpose:}
\begin{itemize}
\item Reduce financial loss due to expired drugs
\item Improve stock rotation
\end{itemize}

\subsection{ ABC Analysis}

Classifies medicines based on economic value contribution.

\textbf{Categories:}
\begin{itemize}
\item A: High value items
\item B: Medium value items
\item C: Low value items
\end{itemize}

\subsection{ EOQ Model (Economic Order Quantity)}

Optimizes order quantity to minimize total inventory cost.

\textbf{Formula:}
\[
EOQ = \sqrt{\frac{2DS}{H}}
\]

Where:
\begin{itemize}
\item D = Demand
\item S = Ordering Cost
\item H = Holding Cost
\end{itemize}

\subsection{ AI Dashboard Simulation}

Provides a consolidated analytical view of the pharmacy system.

\textbf{Outputs:}
\begin{itemize}
\item Total medicines
\item Total sales
\item Revenue
\item Inventory value
\item System recommendations
\end{itemize}

\subsection{Integration Architecture}

\begin{itemize}
\item Node.js handles user interface and workflow control
\item Python executes statistical and simulation models
\item JSON files act as shared data layer
\item Results are printed back to terminal dashboard
\end{itemize}

\subsection{Conclusion}

These simulation models transform the pharmacy system from a basic inventory tool into a predictive analytics platform capable of supporting intelligent pharmaceutical decision-making.


% ---------------- FUTURE ADDITIONS ----------------
\section{Future Additions (Roadmap V13+)}

\begin{tcolorbox}[
colback=lightblue,
colframe=mainblue,
title=System Evolution Roadmap,
arc=3mm,
boxrule=0.8pt]

The Pharmacy Management System is designed as an evolving platform.
Future versions will extend its intelligence, visualization, and clinical utility.

\end{tcolorbox}

\subsection{Advanced Analytics and AI}

\begin{itemize}
\item Predictive demand forecasting for medicines
\item Seasonal inventory optimization
\item AI-based sales trend detection
\item Early warning system for stock depletion
\end{itemize}

\subsection{Visualization and Reporting Enhancements}

\begin{itemize}
\item Embedded charts using \texttt{pgfplots}
\item Revenue trend graphs inside PDF reports
\item Interactive-style dashboard reports (static PDF equivalent)
\item Pharmacy performance scorecards
\end{itemize}

\subsection{Clinical and Pharmacy Intelligence}

\begin{itemize}
\item Drug interaction warning layer
\item Expiry risk prediction system
\item High-risk medicine monitoring
\item Prescription pattern analysis
\end{itemize}

\subsection{System Expansion}

\begin{itemize}
\item Multi-branch pharmacy support
\item Role-based access (Admin / Pharmacist / Cashier)
\item User authentication system
\item Cloud synchronization option
\end{itemize}

\subsection{Automation and Integration}

\begin{itemize}
\item Auto-generated daily/weekly reports
\item Scheduled SAS exports
\item Automatic backup scheduling
\item WhatsApp / SMS notification alerts (future plugin)
\end{itemize}

\subsection{User Experience Improvements}

\begin{itemize}
\item Enhanced CLI interface with structured panels
\item Improved web dashboard UI
\item Mobile-optimized dashboard view
\item Barcode / QR code scanning support
\end{itemize}




% ---------------- CONCLUSION ----------------
\section{Conclusion}

Pharmacy Management System V12 delivers a complete lightweight pharmacy platform combining operational management, analytics, reporting, SAS interoperability, and LaTeX documentation within a single Termux-based ecosystem.





\end{document}




\\section*{Future Additions}

\\begin{itemize}
\\item AI Demand Forecasting
\\item Expiry Prediction
\\item Barcode Scanning
\\item QR Code Support
\\item Multi-Branch Pharmacies
\\item Cloud Synchronization
\\item Advanced SAS Analytics
\\end{itemize}


\\section*{System Highlights}

\\begin{itemize}
\\item Inventory Management
\\item Sales Processing
\\item Analytics Engine
\\item SAS Export
\\item Backup System
\\item Audit Logging
\\item Health Monitoring
\\item Web Dashboard
\\end{itemize}



\\section*{Executive Summary}

This report was generated automatically by the Pharmacy Management System V12.


\\section*{Inventory Summary}

\\begin{longtable}{lrrl}
\\toprule
Medicine & Qty & Price & Expiry \\\\
\\midrule
`;

  m.forEach(x=>{
    tex += `${x.name} & ${x.quantity} & ${x.price} & ${x.expiry} \\\\\n`;
  });

  tex += `
\\bottomrule
\\end{longtable}

\\section*{Sales Summary}

Total Sales: ${s.length}

Total Revenue: ${revenue}

\\begin{longtable}{lrrl}
\\toprule
Medicine & Qty & Total & Date \\\\
\\midrule
`;

  s.forEach(x=>{
    tex += `${x.medicine} & ${x.qty} & ${x.total} & ${x.date} \\\\\n`;
  });

  tex += `
\\bottomrule
\\end{longtable}






